\section{Quality Problems Found at Scale}
\label{sec:quality}

Translating the full corpus revealed two problems that the smaller tests could not show.

\subsection{Characters the Model Cannot Write}
Four strings in the full run came back \emph{empty}: \code{thistle} (Japanese) and \code{Candle}, \code{crab} and
\code{goy Y} (Chinese). Reproducing them with the same model showed the cause: for these strings the model's best output is
the \code{<unk>} (unknown) token, and decoding removes special tokens, leaving nothing. The behaviour is the same at beam~5
and unrelated to batching.

Re-translating the whole Chinese and Japanese sets showed that the empty strings were only the visible part. \textbf{614 of
10,054 Chinese outputs (6.1\%)} and 8 Japanese outputs contained \code{<unk>}. In most of them a character was silently
deleted and the meaning changed: ``New Folder'' became ``new file'', ``Dolphin
Toolbar'' became ``fish tool'', and ``Bitrate'' became just ``rate''.

\subsection{Options Measured}
Three ways of handling these strings were scored against the references on exactly the affected strings
(Table~\ref{tab:unk}).

\begin{table}[H]
\centering
\caption[Handling of \code{<unk>} outputs]{Handling of outputs that contain \code{<unk>} (chrF on the affected strings).}
\label{tab:unk}
\begin{tabularx}{\linewidth}{@{}l S[table-format=2.2] S[table-format=2.2] L@{}}
\toprule
\textbf{Option} & {\textbf{zh (614)}} & {\textbf{ja (8)}} & \textbf{Behaviour} \\
\midrule
Keep the model's text & 26.20 & 6.61 & characters silently missing; meaning often wrong \\
Forbid \code{<unk>} and retry & 25.97 & 33.18 & fixes some; others become ``??'' or a wrong word (``Tunnels'' $\rightarrow$ ``Tunisia'', ``thistle'' $\rightarrow$ ``carbohydrate'') \\
\textbf{Keep the English and report it} & 7.70 & 31.71 & never a wrong word; listed in the log for a translator \\
\bottomrule
\end{tabularx}
\end{table}

\paragraph{Decision.} The pipeline keeps such strings in English and reports them. In localized software, a visibly
untranslated string can be found and fixed; a wrong word cannot. The cost is a lower measured chrF for Chinese, since
chrF counts English as untranslated. Translation Memory entries for these strings are handled in
Chapter~\ref{sec:tmdesign}.

\subsection{Long Strings and the Translation Memory}
The Translation Memory used the English text itself as its database key. SQL Server limits a key to 900 bytes, which
allows at most 434 characters, while KDE4 contains strings of up to 660 characters; saving such a string failed. The key
is now a fixed-size \textbf{SHA-256 hash} of the text, and the text is stored in an unlimited column, so strings of any
length are saved and found again.

\subsection{Verification}
A check run (\code{ec2-20261017t063722z}) translated 383 KDE4 strings (the three longest, 25 known \code{<unk>} cases and
all 355 placeholder strings) in 12 languages, twice, as two identical files. All 24 files succeeded. The second copy was
served \textbf{100\% from the memory} in ten languages, including the long strings; the only misses were the Chinese and
Japanese strings kept in English, which were all listed in the log. Placeholder survival matched Chapter~\ref{sec:placeholders}.

\keyresult{\textbf{RQ4.} At scale, 6.1\% of Chinese outputs silently lost characters the model cannot write, and strings
over 434 characters could not be stored in the memory. Affected strings are now kept in English and reported, and the
memory is keyed by a hash so that strings of any length fit.}
